\subsection{Proof of Reflection}

\label{sec:proof-of-reflection}

In this section we will build up to and discuss a new consensus primitive: \emph{Proof of Reflection} (PoR).
This is a proof that one blockchain's headers have been confirmed by a second chain.
Using these proofs, we can share security between these two chains, such that attacking either chain is as difficult as attacking both.

Our starting point is the question: is there anything \emph{in principle} that prohibits blockchains sharing security?
We already have some examples of this (e.g., merged mining), but those existing methods require each miner to opt-in and aren't general.
Can we come up with a general way to share security between chains?

\aside{
  A note on the term ``miner'': consensus protocols sometimes choose a new name for the role of \emph{block producer}.
  For example: ``validator'', ``baker'', ``collator'', etc.
  In this document, the term ``miner'' usually refers to the generic role of \emph{block producer} in an inclusive sense, rather than specifically to block producers of PoW based-chains.
}

Our starting point is the idea of one chain `tracking' another chain via its headers --- it's the sort of thing that would enable a smart contract to do on-chain SPV proof evaluation.
This isn't a new idea: in 2013, I proposed a system which used this method to support rich cross-chain exchange.\footnote{
  see \url{https://bitcointalk.org/index.php?topic=198032.0}, and
  \url{https://bitcointalk.org/index.php?topic=598784.0}
}
I also wrote a simplified implementation of this method for Bitcoin SPV proofs in the very early days of Ethereum\footnote{
    \url{https://github.com/XertroV/coppr/blob/master/chainheaders.py} \href{https://web.archive.org/web/20240927044057/https://raw.githubusercontent.com/XertroV/coppr/refs/heads/master/chainheaders.py}{[m1]}
} --- a precursor to the later-successful BTC Relay\footnote{
    \url{https://github.com/ethereum/btcrelay}
} (although BTC Relay seems to have been abandoned in late 2017).

What should we call an on-chain version of a headers-only chain?
Since there is no established term, let's call it a ``projection'', and let's call the \emph{act} of one chain creating a projection of another ``imaging''.

\defineTermTex{Projection}{
  A \emph{projection} of a chain is its \emph{headers-only} version that has been recorded and evaluated \emph{by a different chain}.
  For example, \href{https://github.com/ethereum/btcrelay}{BTC Relay} is a smart contract by which Ethereum previously hosted a \emph{projection} of Bitcoin.
  The \emph{act} of one chain creating and maintaining the projection of another is called \emph{imaging}
}

\subsubsection{A Projection of Bitcoin in Ethereum}

It is well understood that an Ethereum smart contract (SC) can \emph{image} Bitcoin.
This results in a \emph{projection} of Bitcoin that is available to other Ethereum SCs (e.g., to use as the basis for SPV proofs).
Since Bitcoin's consensus and PoW algorithm is simple and has low overhead, implementing the necessary logic as an SC is viable.
In principle, any chain with a headers-only mode can be imaged in this way.\footnote{
    Practically, there can be constraints that prevent this.
    For example: Ethereum typically doesn't support memory hard hashes, so hosting projections of chains like Litecoin on Ethereum is potentially impossible.
}

Let's assume that a smart contract like this exists on the Ethereum chain and it is operational.

We're going to use some tables and diagrams as we go to show the sequence of events that occur and what data is produced and/or recorded.
\autoref{tab:eth-images-btc} shows data and events for both Bitcoin and Ethereum as Bitcoin blocks are produced and the projection of Bitcoin in Ethereum is maintained.
\autoref{fig:pr-btc-eth-step1} illustrates this.
Note: \autoref{fig:pr-btc-eth-step1} includes some variance in Ethereum's block production rate, similar to what might be observed in a real-world environment, but \autoref{tab:eth-images-btc} does not.

After a Bitcoin block ($\text{BTC}_k$) is produced, a user (it doesn't matter who) produces a transaction informing the SC of that block's header.
Then, an Ethereum miner includes that transaction in a block, which updates the SC.
This process repeats when the next Bitcoin block is mined, and so on.

Why would a chain want\footnote{
    Obviously, blockchains don't have desires.
    We'll personify them as a short-cut, which here means approximately: the users / community / developers of a blockchain have some common goal or incentive to do a particular thing.
} to include a projection of another chain?
The typical answer is to enable proofs of some state or that some transactions occurred on the imaged chain.
This enables more complex SCs, e.g., a trustless $\text{BTC}\leftrightarrow\text{ETH}$ market.


\ctable[
  pos = h,
  caption = (Hypothetical) Data and events for both Bitcoin and Ethereum as a projection of Bitcoin in Ethereum is maintained via Bitcoin headers being tracked by an Ethereum SC.,
  cap = (Hypothetical) Data and events as Ethereum images Bitcoin.,
  center,
  label = tab:eth-images-btc,
]{lllll}{}{
  \FL
  \shortstack[l]{Time step (\textasciitilde{}15 s \\ increments)} & \shortstack[l]{Bitcoin \\ block mined} & \shortstack[l]{Eth \\ block mined} & \shortstack[l]{Eth block contents} & {Eth state}
  \ML
  $\vdots$ & & & & \NN
  0 & $k$ & & & \NN
  1 & & $j$ & $\text{BTC}_k$ header & Records $\text{BTC}_{0 \cdots k}$ \NN
  $\vdots$ & & & & \NN
  40 & $k+1$ & & & \NN
  41 & & $j+40$ & $\text{BTC}_{k+1}$ header & Records $\text{BTC}_{0 \cdots k+1}$ \NN
  $\vdots$ & & & &
  \LL
}

\begin{figure}
\centering
\includegraphics[max width=\linewidth, height=0.35\textheight]{pow_refl_btc_eth_step1_sag}
\caption[
  (Hypothetical) A projection of Bitcoin in Ethereum via an SC and transactions.
]{
  (Hypothetical) Bitcoin headers, as they are produced, are included in Ethereum's state via a smart contract and user made transactions.
  This is roughly how BTC Relay worked.
}
\label{fig:pr-btc-eth-step1}
\end{figure}


\subsubsection{Incrementally Implementing \emph{Proof of Reflection}}

\label{sec:two-blockchains}

Let's consider two hypothetical and distinct blockchains --- chain \cL and chain \cR.
To begin with, you can imagine these as Bitcoin and Ethereum if it helps.
However, the changes that are required by \emph{Proof of Reflection} would require major changes to either's consensus mechanism, so it's unlikely that changes like these would ever be integrated.

Our starting case for \cL and \cR is that both use Nakamoto consensus with the same Proof of Work algorithm.
For simplicity, both chains also have identical block times and we won't consider block production variance.

\subsubsubsection{Step 1. Chain \cR Images Chain \cL}

This is similar to Ethereum imaging Bitcoin --- see \autoref{fig:pow_refl_step1}.
(We'll omit the table this time since it's basically the same as \autoref{tab:eth-images-btc}.)

Like before, \cR will include \cL's headers as they are produced.
Unlike before, let's assume this is a \emph{protocol-level} implementation; so \cR miners can include \cL headers directly with no transaction fees.

\begin{figure}
    \centering
    \includegraphics[max width=\linewidth, height=0.23\textheight]{pow_refl_step1_sag}
    \caption{Step 1. Chain R images Chain L; thus Chain R hosts a \emph{projection} of Chain L.}
    \label{fig:pow_refl_step1}
\end{figure}

\subsubsubsection{Step 2. \cL Images \cR}

Similar to \cR, \cL will begin including \cR headers in each block via appropriate protocol-level changes.
Just as \cR hosts a projection of \cL, \cL now hosts a projection of \cR.
This is shown in \autoref{fig:pow_refl_step2} and \autoref{tab:por-step-2}.

\begin{table}[hb]
\centering
\caption{Step 2. Both Chain L and Chain R host a projection of each other.}
\label{tab:por-step-2}
\resizebox{\textwidth}{!} {%
\begin{tabular}{lllllll}
\toprule
{Time} & \shortstack[l]{L block \\ made} & \shortstack[l]{L block \\ contents} & {L state} & \shortstack[l]{R block \\ made} & \shortstack[l]{R block \\ contents} & {R state} \\
\midrule
{$\vdots$} & {} & {} & {} & {} & {} & {} \\
{0} & $k$ & {$R_{j-1}$ header} & {Records $R_{0 \cdots j-1}$} & {} & {} & {} \\
{1} & {} & {} & {} & $j$ & {$L_{k}$ header} & {Records $L_{0 \cdots k}$} \\
{2} & $k + 1$ & {$R_{j}$ header} & {Records $R_{0 \cdots j}$} & {} & {} & {} \\
{3} & {} & {} & {} & $j + 1$ & {$L_{k+1}$ header} & {Records $L_{0 \cdots k+1}$} \\
{$\vdots$} & {} & {} & {} & {} & {} & {} \\
\bottomrule
\end{tabular}%
}
\end{table}

\begin{figure}
  \centering
  \includegraphics[max width=\linewidth, max height=0.3\textheight]{pow_refl_step2_sag}
  \caption{Step 2. Chain L and Chain R contain a projection of each other's headers-only chain.}
  \label{fig:pow_refl_step2}
\end{figure}

\FloatBarrier

% The following was the beginning of a redraft, and the content is still in 20-proof-of-reflection.md

% \subsubsubsection{Step 3. A Reflection of \cL in \cR}

% Can we use a projection of a chain for a new purpose?
% What might be possible if \cL keeps track of \cR's projection of \cL?
% This is straightforward to do if \cL and \cR expose their projections via a commitment in the header.\footnote{
%     This is commonly done via the root of a merkle or verkle tree, where that tree contains state entries of some kind.
% }

% \begin{figure}
%   \centering
%   \includegraphics[max width=\linewidth, height=0.28\textheight]{pow_refl_step3_sag}
%   \caption{Step 3 (todo)}
%   \label{fig:pow_refl_step3}
% \end{figure}
